The literature on automated AI-text detection spans roughly two camps: \emph{statistical/unsupervised}
methods that require no labelled data, and \emph{supervised} classifiers trained on paired
human-versus-machine corpora. We briefly survey each and situate RADAR within this landscape.

\subsection{Zero-Shot and Statistical Detection}

The earliest approaches to automatic authorship verification treated the problem as a
one-class classification task~\cite{koppel:2004}. Statistical compression-based methods
model a text's regularity against a reference model: PPMd compression-based cosine
similarity~\cite{sculley:2006,halvani:2017} computes how predictably a corpus model
can encode an unknown document. These methods are inherently robust to topic drift
but suffer on short texts and on cleverly paraphrased outputs.

Perplexity-based detectors leverage the observation that LLMs assign higher probability
to their own outputs than to human-written text. DetectGPT~\cite{mitchell:2023} exploits
local probability curvature: AI-generated text tends to lie near local maxima of the model's
probability surface. Fast-DetectGPT~\cite{bao:2023} accelerates this by using conditional
probability curvature rather than explicit perturbation sampling. DetectLLM~\cite{su:2023}
uses log-rank information from a surrogate language model. Binoculars~\cite{hans:2024}
achieves strong zero-shot performance by cross-referencing perplexity and cross-perplexity
from two related models, reducing the impact of topic-specific language complexity.

A critical limitation of all purely statistical approaches is their dependence on a
reference LM that should match the generator. When the generator is unknown—as in
competitive shared tasks—performance degrades significantly~\cite{sadasivan:2023}.
Moreover, adversarial attacks such as paraphrase and synonym substitution substantially
reduce statistical signals because they alter the surface-level probability fingerprint
while preserving the underlying content~\cite{ippolito:2020}.

\subsection{Supervised Classification}

Supervised approaches train discriminative models on labelled datasets. The linear SVM
with TF-IDF features~\cite{bevendorff:2025} provided as the task baseline demonstrates
that even simple bag-of-words models are surprisingly effective on curated in-domain
corpora. However, they generalise poorly to out-of-distribution generators and to
adversarially obfuscated texts.

Fine-tuned transformer models have emerged as strong baselines. BERT- and RoBERTa-based
classifiers achieve high accuracy on standard benchmarks such as
TuringBench~\cite{uchendu:2021} and the datasets from the PAN 2024
campaign~\cite{bevendorff:2024b}. The introduction of larger encoders—GPT2-based detectors,
LLaMA fine-tunes—has extended performance further, but these models are prone to
distribution shift when the test set comes from a different generator family than the
training data~\cite{schuster:2020,sadasivan:2023}.

A key insight from recent competitive results is that pure end-to-end neural approaches
are not unambiguously superior to hybrid models that combine neural representations
with handcrafted features. Stylometric features—vocabulary richness, syntactic
complexity, function word distributions—tend to be more stable across paraphrase
perturbations because they reflect deep properties of the generation process rather
than surface token choices~\cite{bevendorff:2025b,uchendu:2020}.

\subsection{Stylometric Analysis in AI Detection}

Stylometry has a long history in authorship attribution and verification~\cite{koppel:2004,
bevendorff:2019}. For AI detection, the stylometric hypothesis is that LLMs, despite
persona-mimicry prompting, exhibit systematic statistical regularities that differ from
human writing: higher lexical homogeneity, more regular sentence length distributions,
lower pronoun density, and more uniform paragraph coherence~\cite{jakesch:2023,
schuster:2020}. These properties persist under paraphrase because they reflect the
model's training-time inductive biases rather than its surface vocabulary.

\subsection{Multi-Source Training and Robustness}

The PAN 2024 overview~\cite{bevendorff:2024d} identified domain shift and adversarial
obfuscation as the two primary failure modes of current detectors. RAID~\cite{dugan:2024}
was introduced specifically to study this: a benchmark of 5.6~million samples from 11~LLMs
across 11~domains, with 12 adversarial attack types including paraphrase, synonym substitution,
and homoglyph character substitution. Training on RAID has been shown to substantially
improve robustness, but only when the training distribution is appropriately balanced.

\subsection{Positioning RADAR}

RADAR is designed to directly address the key challenges identified above. By fusing
deep semantic representations from a state-of-the-art encoder with handcrafted
stylometric features through a cross-attention mechanism, RADAR captures both
surface-level patterns and deeper structural regularities. The two-stage training
curriculum—frozen encoder followed by full fine-tuning on a balanced multi-source
corpus—provides strong initialisation before allowing end-to-end adaptation. This
stands in contrast to single-stage approaches that either sacrifice stylometric
signals (pure neural) or discard contextual reasoning (pure stylometric).
